% ============================================================
%  Глава 13 — Геозона, преобразование сенсоров, реактивная навигация
%  Файлы: geofence_map_publisher.py, depth_to_scan_node.py, fallback_nav_node.py
% ============================================================
\chapter{Геозона, преобразование сенсоров и реактивная навигация}
\label{ch:geofence_sensors}

% ═══════════════════════════════════════════════════════════════
\section{Геозона: алгоритм Point-in-Polygon}

\begin{filebox}[title={Файл исходного кода}]
\filepath{compute\_node/geofence\_map\_publisher.py} \quad (108 строк, Python)
\end{filebox}

\begin{filebox}[title={\filepath{geofence\_map\_publisher.py}, строки 79--91 --- \_point\_in\_polygon()}]
\begin{lstlisting}[style=pythonstyle, numbers=left, firstnumber=79]
@staticmethod
def _point_in_polygon(x, y, poly):
    """Ray-casting algorithm."""
    n = len(poly)
    inside = False
    j = n - 1
    for i in range(n):
        xi, yi = poly[i]
        xj, yj = poly[j]
        if ((yi > y) != (yj > y)) and \
           (x < (xj - xi) * (y - yi) / (yj - yi) + xi):
            inside = not inside
        j = i
    return inside
\end{lstlisting}
\end{filebox}

\subsection{Алгоритм трассировки луча (Ray Casting)}

Для определения, находится ли точка $(x, y)$ внутри полигона,
из неё проводится горизонтальный луч вправо и считается число
пересечений с рёбрами:

\begin{equation}
  \text{inside} = \bigl(\text{число пересечений}\bigr) \bmod 2 = 1
  \label{eq:raycasting}
\end{equation}

Для каждого ребра $(x_i, y_i) \to (x_j, y_j)$:

\textbf{1. Проверка пересечения по $y$:}
\begin{equation}
  (y_i > y) \neq (y_j > y) \quad\text{(ребро пересекает горизонталь $y$)}
  \label{eq:edge_cross}
\end{equation}

\textbf{2. Вычисление $x$-координаты пересечения} (линейная интерполяция):
\begin{equation}
  x_{\text{cross}} = x_i + \frac{(x_j - x_i)(y - y_i)}{y_j - y_i}
  \label{eq:x_intersect}
\end{equation}

\textbf{3. Если} $x < x_{\text{cross}}$ --- луч пересекает ребро, инвертируем флаг.

\subsection{Построение карты занятости}

\begin{filebox}[title={\filepath{geofence\_map\_publisher.py}, строки 46--77}]
\begin{lstlisting}[style=pythonstyle, numbers=left, firstnumber=47]
poly = np.array(self._polygon, dtype=np.float64)
x_min = poly[:, 0].min() - self._padding   # 0.5 m padding
x_max = poly[:, 0].max() + self._padding
width  = int((x_max - x_min) / self._resolution)  # 0.05 m/cell
height = int((y_max - y_min) / self._resolution)

# 0=free inside, 100=lethal outside
grid = np.full((height, width), 100, dtype=np.int8)
for iy in range(height):
    for ix in range(width):
        px = x_min + ix * self._resolution
        py = y_min + iy * self._resolution
        if self._point_in_polygon(px, py, poly):
            grid[iy, ix] = 0
\end{lstlisting}
\end{filebox}

Карта создаётся как \texttt{OccupancyGrid} (ROS2 nav\_msgs):
\begin{itemize}
  \item $0$ --- свободная ячейка (внутри полигона)
  \item $100$ --- летальное препятствие (за пределами арены)
\end{itemize}

% ═══════════════════════════════════════════════════════════════
\section{Ультразвук $\to$ LaserScan: синтетическое сканирование}

\begin{filebox}[title={Файл исходного кода}]
\filepath{compute\_node/depth\_to\_scan\_node.py} \quad (82 строки, Python)
\end{filebox}

\begin{filebox}[title={\filepath{depth\_to\_scan\_node.py}, строки 43--65}]
\begin{lstlisting}[style=pythonstyle, numbers=left, firstnumber=43]
def _publish_scan(self):
    scan = LaserScan()
    scan.angle_min = -math.pi / 2.0        # -90 deg
    scan.angle_max =  math.pi / 2.0        # +90 deg
    scan.angle_increment = (ANGLE_MAX - ANGLE_MIN) / NUM_BEAMS  # 180 beams
    scan.range_min = 0.02
    scan.range_max = 3.0

    ranges = [MAX_RANGE] * NUM_BEAMS        # all unknown
    centre = NUM_BEAMS // 2
    # Spread ultrasonic over ~15 deg cone
    cone_beams = int(0.26 / scan.angle_increment)
    for i in range(max(0, centre - cone_beams),
                   min(NUM_BEAMS, centre + cone_beams + 1)):
        ranges[i] = min(self._latest_range, MAX_RANGE)
    scan.ranges = ranges
\end{lstlisting}
\end{filebox}

\subsection{Геометрия конуса ультразвука}

Один ультразвуковой датчик HC-SR04 даёт единственное расстояние
с конусом обзора ${\approx}30^\circ$. Для совместимости с slam\_toolbox
это значение «размазывается» по нескольким лучам:

\begin{equation}
  \Delta\alpha = \frac{\alpha_{\max} - \alpha_{\min}}{N_{\text{beams}}}
  = \frac{\pi}{180} \approx 0.0175\,\text{рад/луч}
  \label{eq:angle_incr}
\end{equation}

\begin{equation}
  N_{\text{cone}} = \left\lfloor \frac{0.26}{\Delta\alpha} \right\rfloor
  \approx 15\,\text{лучей в полуконусе}
  \label{eq:cone_beams}
\end{equation}

Итого ${\approx}30$ центральных лучей получают реальное измерение,
остальные 150 --- значение $r_{\max} = 3.0\,\text{м}$ (неизвестно).

% ═══════════════════════════════════════════════════════════════
\section{Реактивная навигация (fallback)}

\begin{filebox}[title={Файл исходного кода}]
\filepath{pi\_nodes/nodes/fallback\_nav\_node.py} \quad (176 строк, Python)
\end{filebox}

\begin{filebox}[title={\filepath{fallback\_nav\_node.py}, строки 122--150 --- \_run\_fallback\_drive()}]
\begin{lstlisting}[style=pythonstyle, numbers=left, firstnumber=136]
if r > SAFE_M:                          # r > 0.40 m
    cmd['linear_x'] = FWD_SPEED        # 0.10 m/s
    self._state = STATE_FORWARD
elif r > STOP_M:                        # 0.20 < r < 0.40
    scale = (r - STOP_M) / (SAFE_M - STOP_M)
    cmd['linear_x'] = SLOW_SPEED + scale * (FWD_SPEED - SLOW_SPEED)
    self._state = STATE_CAUTION
else:                                   # r < 0.20 m
    self._state = STATE_ROTATE
    cmd['angular_z'] = ROT_SPEED * self._rotate_dir  # +/- 0.8 rad/s
\end{lstlisting}
\end{filebox}

\subsection{Пропорциональное масштабирование скорости}

В зоне осторожности ($r_{\text{stop}} < r < r_{\text{safe}}$):

\begin{equation}
  \text{scale} = \frac{r - r_{\text{stop}}}{r_{\text{safe}} - r_{\text{stop}}}
  = \frac{r - 0.20}{0.40 - 0.20} \in [0, 1]
  \label{eq:caution_scale}
\end{equation}

\begin{equation}
  \boxed{v = v_{\text{slow}} + \text{scale} \cdot (v_{\text{fwd}} - v_{\text{slow}})}
  \label{eq:caution_speed}
\end{equation}

Это \textbf{линейная интерполяция} между $v_{\text{slow}} = 0.05\,\text{м/с}$
и $v_{\text{fwd}} = 0.10\,\text{м/с}$.

\subsection{Конечный автомат реактивной навигации}

\begin{center}
\begin{tikzpicture}[
  state/.style={rectangle, rounded corners, draw=blue!60, fill=blue!10,
                minimum width=2.5cm, minimum height=0.9cm, font=\small\bfseries},
  arrow/.style={-{Stealth[length=6pt]}, thick}
]
  \node[state] (F) {FORWARD};
  \node[state, right=3cm of F] (C) {CAUTION};
  \node[state, below=1.5cm of C] (R) {ROTATE};

  \draw[arrow] (F) to node[above, font=\tiny]{$r < 0.40$} (C);
  \draw[arrow] (C) to node[right, font=\tiny]{$r < 0.20$} (R);
  \draw[arrow] (R) to[bend right=30] node[left, font=\tiny]{$r > 0.40$} (F);
  \draw[arrow] (C) to[bend left=20] node[above, font=\tiny]{$r > 0.40$} (F);
\end{tikzpicture}
\end{center}

\subsection{Теоретическое обоснование: конечный автомат как дискретная система}

Реактивная навигация реализует \textbf{гибридную систему} ---
комбинацию непрерывной динамики (скорость) и дискретных переключений
(состояния автомата). В терминах ТАУ это \textbf{переключающаяся система}:

\begin{equation}
  \dot{\vect{x}} = \vect{f}_{q(t)}(\vect{x}), \quad q \in \{\text{FWD}, \text{CAUT}, \text{ROT}\}
  \label{eq:switching_system}
\end{equation}

\textbf{Устойчивость} переключающейся системы гарантирована, если
каждая подсистема устойчива и переключения не слишком частые
(условие dwell time).

Линейная интерполяция~\eqref{eq:caution_speed} в зоне CAUTION
обеспечивает \textbf{непрерывность} скорости на границах переключений,
предотвращая рывки.

\section{Сводная таблица параметров}

\begin{center}
\begin{tabular}{llll}
\toprule
\textbf{Модуль} & \textbf{Параметр} & \textbf{Значение} & \textbf{Описание}\\
\midrule
Геозона & Разрешение & $0.05\,\text{м}$ & Размер ячейки\\
Геозона & Отступ & $0.5\,\text{м}$ & Запас за полигоном\\
LaserScan & $N_{\text{beams}}$ & $180$ & Кол-во лучей\\
LaserScan & FOV & $180^\circ$ & Поле зрения\\
Fallback & $r_{\text{safe}}$ & $0.40\,\text{м}$ & Безопасная дистанция\\
Fallback & $r_{\text{stop}}$ & $0.20\,\text{м}$ & Стоп-дистанция\\
Fallback & $\omega_{\text{rot}}$ & $0.8\,\text{рад/с}$ & Скорость поворота\\
\bottomrule
\end{tabular}
\end{center}
